\documentclass[12pt]{article}
\usepackage[utf8]{inputenc}
\usepackage[T1]{fontenc}
\usepackage{hyperref}
\title{Zan Gula Language Resources for K-12 Curriculum Development}
\author{Compiled by Claude Sonnet 5.5 (claude-sonnet-5-5) for Andr\'as Kornai}
\date{}
\begin{document}
\maketitle

\section{Introduction}
Zan Gula (ISO 639-3 zna, Glottocode zang1249) is a language of Chad. Wikipedia describes it as an Adamawa language. Alternate names include Gula and Moriil. Speaker estimates range from about 5,000 to 15,000 depending on the source. Search results turned up very little material for this language.

\section{Primary Educational Resources}
No primer, textbook or curriculum was found in search results. Although a reference description says the language is taught as a subject in some schools, no source for that turned up.

\section{Linguistic and Reference Materials}
\begin{itemize}
\item Wikipedia article on Zan Gula: \url{https://en.wikipedia.org/wiki/Zan_Gula_language}
\item Numeral system dataset for Gula, Moyen-Chari region, Chad, from the Max Planck Institute (it may describe a related variety or a broader grouping): \url{https://lingweb.eva.mpg.de/channumerals/Gula-Chad.htm}
\end{itemize}

\section{Religious and Community Resources}
\begin{itemize}
\item Joshua Project profile of Zan Gula in Chad: \url{https://joshuaproject.net/people_groups/16084}
\item Joshua Project language page for zna: \url{https://joshuaproject.net/index.php/languages/zna}
\end{itemize}

\section{Digital and Audio Resources}
Global Recordings Network has a language page for zna: \url{https://globalrecordings.net/en/language/zna}. The search results did not show which recordings it lists.

\section{Curriculum Development Context}
The documentation found is minimal. Any curriculum would require new fieldwork, an agreed orthography and community partnership. French and Chadian Arabic are the languages of wider communication in Chad. The sources found here are background references, not teaching materials.

\end{document}